\subsection{Technical Insights and Trade-offs}

These clinical observations reflect underlying technical trade-offs in modality selection, architecture design, and personalization strategy that shape what is achievable with current wearable approaches.

\subsubsection{Modality Paradox}

The evidence reveals a paradox regarding modality selection. Multi-modal approaches dominate the research field, with 69\% of studies using multiple sensor types. However, the study achieving the lowest FPR uses single-modality accelerometry. This finding suggests that sensor fusion is not always necessary for reliable detection, and that biological specificity of motor manifestations may be more important than modality count.

The physiological requirements differ between detection and forecasting. Detection can succeed with movement sensors alone because ictal motor manifestations produce clear signals in ACC recordings. Forecasting requires autonomic modalities because pre-ictal changes must be detected before clinical symptoms appear. All forecasting studies incorporate EDA, ECG, or temperature data to capture these autonomic precursors. The contrast suggests that modality selection should be guided by the temporal phase of seizure activity being targeted.

\subsubsection{Architecture Selection}

Detection and forecasting applications employ different architectural patterns. Detection research shows methodological diversity including CNNs, LSTMs, ANNs, ensemble methods, and anomaly detection. This diversity may reflect the broader exploration phase of detection research. Forecasting studies concentrate heavily on LSTM variants, with three of four studies using this architecture. This convergence reflects the different temporal demands, as forecasting requires modeling longer-term dependencies than detection.

The persistence of competitive handcrafted feature approaches is noteworthy. Domain knowledge remains valuable despite the end-to-end learning trend, with engineered features achieving excellent performance in some studies. This suggests that purely data-driven approaches may not fully grasp the known physiological signatures of seizures or that current end-to-end approaches can be optimized further.

Ensemble methods achieve excellent results but require greater computational resources. The trade-off between ensemble complexity and deployability becomes important for commercial translation, as multi-model ensembles may be impractical for low-power wearable devices.

\subsubsection{Personalization Dilemma}

The personalization strategy presents a fundamental dilemma for clinical deployment. Global models enable immediate deployment without individual patient data, but only 43\% of patients achieve performance above chance in LOSO validation. This means most patients receive no meaningful protection from population-based systems rendering them useless for clinical deployment.

Patient-specific approaches achieve 100\% success rates but require substantial individual data collection. The cold-start problem means newly diagnosed patients must wait weeks or months before the system becomes useful. Mixed approaches combining population-level patterns with individual data may offer a middle path, but evidence remains limited to single studies.

\subsubsection{The False Positive Challenge}

The false positive challenge differs substantially between modalities. Accelerometer-based approaches can achieve clinically acceptable FPR below 0.05/h, but ECG-based approaches produce FPR an order of magnitude higher. Even the best ECG-based method exceeded the clinical benchmark by a substantial margin. This suggests that cardiac signals alone cannot reliably distinguish seizures from normal daily activities.

Multi-modal sensor fusion does not automatically solve the false positive problem. Studies combining multiple modalities still report FPR above clinical targets. The challenge may stem from the fundamental similarity between seizure-related physiological changes and those occurring during normal activities such as exercise, stress, or sleep transitions.

Forecasting faces a different performance ceiling. The consistency of AUC around 0.75 across different methods and modalities suggests fundamental limits on pre-ictal predictability using non-invasive sensors. This ceiling may reflect genuine biological variability in pre-ictal states that cannot be resolved with current wearable modalities and requires furhter research on the exact biological mechanisms during an epileptic seizure.
